\section{Questions}
\begin{enumerate}
    \item Notre conception actuelle utilise un cône arrière (\textit{tail cone}) composé de deux parties : (1) un cône arrière Pro54 du fabricant, et (2) une surface supplémentaire placée dessus afin qu'il soit au ras du tube de la fusée. Nous comprenons que les cônes arrière ne sont pas fournis par le C'Space. Nous en achèterons un pour une autre compétition, mais nous aimerions savoir si nous pouvons utiliser cette conception au C'Space, ou si elle n'est pas autorisée.

    \item Nous croyons nous souvenir qu'au C'Space, il est interdit de larguer quoi que ce soit depuis la fusée (sauf si l'élément reste attaché), pour des raisons de sécurité et de protection de l'environnement. Cette année, nous prévoyons de développer un bouton de rail qui se détache pendant le lancement. Cela serait-il acceptable ? Notre conception s'inspire de cette vidéo : \url{https://youtu.be/LnunMEopsb8?t=83}

    \item If we are allocated the Pro75 motor, would we be able to switch to the Pro54 if a design change comes up later? This is unlikely, but we would like to know just in case.

    \item What are the rail dimensions? We have only been able to find the maximum rocket dimensions, not the rail size, which would be useful for the rocket's vehicle dynamics.

    \item Regarding the 3000\,m apogee limit, we understand it is a maximum. Should we aim slightly lower, e.g.\ within 100\,m of the 3000\,m limit? Or is it acceptable as long as StabTraj predicts less than 3000\,m?

    \item If we fly transonic, would it be acceptable to still use composite fins provided we demonstrate their strength? We are in contact with another team using the same method we intend to use; their predicted maximum loading is 35\,kg per fin, which should be sufficient.
\end{enumerate}